\section{考察}
\subsection{検出器の選択}
検出の精度は，評価する物体によって順位が異なった．
COCOの物体に限れば，YOLO11xとDETRは少ないボックスで物体を囲み，GTに当たらないボックスも少ない．
そのため，COCOの80クラスで評価した検出から命名までの$ F $は，YOLO11xのボックスを用いた構成が最も高かった．
一方，COCO外の物体の再現率は，YOLO11xが0.100，DETRが0.124にとどまり，閾値を下げても0.18以下であった．
COCOで学習した検出器は，学習したクラスの物体以外を背景として扱うよう学習されているためである．

本稿の目的は，再学習なしに，COCOの80クラスより細かい語や80クラスの外の語で物体をアノテーションすることである．
検出器が囲まない物体は，分類器や語彙の木をどれだけ改善してもアノテーションできない．
したがって，検出器の再現率はアノテーションできる物体の上限を決める．
OWLv2の物体らしさは，学習時のクラスに依存せず，COCOの物体の87.5\%とCOCO外の物体の73.4\%を囲む唯一の検出器であった．
また，ボックス数$ B $により再現率を連続的に調整できる．
本稿でOWLv2を検出器に選ぶ理由はこの2点である．
YOLO11xやDETRを用いると，COCOの物体の$ F $は高まるが，COCO外の物体のアノテーションを諦めることになる．
なお，COCOの物体ではYOLO11xのボックスが有利であるため，YOLO11xのボックスとOWLv2のボックスを組み合わせることも有望である．

\subsection{分類器の選択}
Alpha-CLIPは，GTのボックスのTop-1と部分木の確率の合計による命名の$ hF $でCLIPを上回り，cropの入力では全ての検出器との組み合わせで$ F $がCLIPを上回った．
Alpha-CLIPは画像全体の文脈を保ったまま，αチャネルで指定した領域の特徴を求めるため，ボックスの内側だけを切り抜くCLIPより物体の識別に有利である．
また，COCOのクラスより細かい語を正しく出力した割合はCLIPの約2倍であり，粒度を変更するモデルの分類器として適している．
一方，命名の正解率はCLIPの方が高く，Alpha-CLIPは細かい語を出力する代わりに誤りも増える．
細かい語と正しさのどちらを優先するかは，式(\ref{eq_output})の閾値$ \tau $で調整できる．
contextの入力は周囲の物体の名前が付きやすく，ボックスの多いOWLv2との組み合わせではcropの入力が適していた．

\subsection{語彙の木による粒度の判断}
WordNetの下位語を1段ずつたどる命名では，personの直接の下位語が398語あり，1段目で正しい子を選べる割合は0.005にとどまった．
\rev{WordNet$ \cap $VGの木は，VGの注釈に現れる語のみで木を作るため，}personの子を136語に絞り，1段目の正解率を0.097に高めた．
さらに，部分木の確率の合計は，1段ずつの選択の誤りを積み重ねず，全ての語の確率から出力の深さを決める．
VLMの確率が人物の多くの下位語に分散する場合でも，部分木の合計はpersonのノードに集まるため，細かい語に確信が無ければpersonを出力する．
これが，GTのボックスで$ hF $が0.266から0.639に向上した理由である．
ただし，GTのボックスで正しくCOCOのクラスより細かい語を出力した割合は最大でも0.123であり，細かい粒度の命名の正しさには改善の余地がある．

\subsection{本稿の限界と今後の課題}
第一に，主な評価はCOCOの80クラスと，85\%が人物であるキャプションの細かい正解に基づく．
この評価はCOCOで学習した検出器に有利であり，COCO外の物体の命名の正しさは評価していない．
LVISのカテゴリはWordNetのSynsetを持つため，LVISの注釈を細かい正解とすれば，人物以外の物体で深さごとの命名を評価できる．
第二に，LVISの注釈は網羅的でなく，COCO外の物体の再現率と誤検出の内訳は概算である．
第三に，細かい正解はキャプションに基づくため，個々のインスタンスとは対応せず，第1語義の誤りを含む．
第四に，$ B = 50 $のOWLv2の推論のため，処理時間はOldの約5倍である．
第五に，残すボックスを決める規則は，COCOの80クラスの部分木の確率の合計と物体らしさの積を用い，$ \theta $をCOCOの$ F $で選ぶため，COCO外の物体のボックスは残らない．
その結果，検出器の出力では0.734であったCOCO外の物体の再現率が，残したボックスでは0.103に低下した．
COCO外の物体を残す棄却のスコアの設計と，その評価は今後の課題である．
